%% Regression: whole-cell overrides suppress basis members before spacing.
%% A partially overridden window still diagnoses realized collisions, while a
%% collision that requires a suppressed member remains quiet.  A
%% member-addressed replacement remains a live member, and a spanning member
%% contributes its canonical anchor rather than aliases at every claimed cell.
%% A suppressed picture checks that no live-candidate state leaks.  An
%% out-of-window span defeats count-only completeness without inventing a
%% candidate, while a translated inside member can still collide with that
%% outside anchor.  The final dense pictures trap sparse-product construction;
%% an additional half-pitch case traps occupancy scans before a raw violation.
\documentclass{standalone}
\usepackage{tenkz}
\makeatletter
\ExplSyntaxOn
\file_input:n { tenkz-kernel.code.tex }
\ExplSyntaxOff
\makeatother
\begin{document}
\tenkzkernel
\begin{tenkz}[rows={ket}, cols=1, bonds=none,
              frame={flat,basis={wire at (0,0), wire at (0,0)}}]
  \tn[at=(1,1), name=X]{X}
  \tnmark[form=enclosure]{(1,1)}{}
\end{tenkz}
\begin{tenkz}[rows={ket}, cols=2, bonds=none,
              frame={flat,basis={wire at (0,0), wire at (0,0)}}]
  \tn[at=(1,1), name=Y]{Y}
\end{tenkz}
\begin{tenkz}[rows={ket}, cols=2, bonds=none,
              frame={flat,basis={wire at (0,0), wire at (4,0)}}]
  \tn[at=(1,1), name=Z]{Z}
\end{tenkz}
\begin{tenkz}[rows={ket}, cols=1, bonds=none,
              frame={flat,basis={wire at (0,0), wire at (0,0)}}]
  \tn[at=(1,1,1), name=M]{M}
\end{tenkz}
\begin{tenkz}[rows={ket}, cols=2, bonds=none,
              frame={flat,basis={wire at (0,0), wire at (4,0)}}]
  \tn[at=(1,1,1), wide=2, name=W]{W}
\end{tenkz}
\begin{tenkz}[rows={ket}, cols=2, bonds=none,
              frame={flat,basis={wire at (4,0), wire at (0,0)}}]
  \tn[at=(1,1,1), wide=2, name=V]{V}
\end{tenkz}
\begin{tenkz}[rows={ket}, cols=1, bonds=none,
              frame={flat,basis={wire at (0,0), wire at (0,0)}}]
  \tn[at=(1,1), name=Q]{Q}
\end{tenkz}
\begin{tenkz}[rows={ket}, cols=1, bonds=none,
              frame={flat,basis={wire at (0,0), wire at (0,0)}}]
  \tn[at=(1,0,1), wide=2, name=O]{O}
\end{tenkz}
\begin{tenkz}[rows={ket}, cols=2, bonds=none,
              frame={flat,basis={wire at (0,0), wire at (-4,0)}}]
  \tn[at=(1,0,1), wide=2, name=P]{P}
\end{tenkz}
\begin{tenkz}[rows={ket, ket}, cols=2, bonds=none,
              frame={flat,basis={wire at (0,0), wire at (0,0)}}]
\end{tenkz}
\ExplSyntaxOn
\cs_set_protected:Npn \__tenkz_kernel_basis_spacing_sparse_outer:n #1
  { \msg_fatal:nnn {tenkz}{kernel-frame-basis-parse}{dense-plan} }
\ExplSyntaxOff
\begin{tenkz}[rows={ket, ket, ket}, cols=3, bonds=none,
              frame={flat,basis={wire at (0,0), wire at (0,0)}}]
  \tn[at=(3,3), name=H]{H}
\end{tenkz}
\ExplSyntaxOn
\cs_set_protected:Npn
  \__tenkz_kernel_basis_spacing_dense_realized:nnnn #1#2#3#4
  { \msg_fatal:nnn {tenkz}{kernel-frame-basis-parse}{nonnegative-scan} }
\ExplSyntaxOff
\begin{tenkz}[rows={ket, ket, ket}, cols=3, bonds=none,
              frame={flat,basis={wire at (0,0), wire at (2,0)}}]
  \tn[at=(3,3), name=S]{S}
\end{tenkz}
\end{document}
